\documentclass[10pt,twocolumn]{article}
\usepackage[margin=0.75in]{geometry}
\usepackage{amsmath,amssymb,booktabs,xcolor}
\usepackage[hidelinks]{hyperref}
\usepackage{caption}
\captionsetup{font=small,labelfont=bf}
\newcommand{\JobID}{dc0uhdts2omc7389tmt0}
\newcommand{\QPUs}{63}
\newcommand{\QPYsha}{d78d15b7a2040fb1}
\newcommand{\NCirc}{216}
\newcommand{\QSha}{415fb1bd32be1e51}
\newcommand{\HQCtotalC}{8,754}
\newcommand{\HQCtotalA}{2,039}
\newcommand{\HQCtotalB}{4,557}
\newcommand{\HWTable}{1 & 2 & 0.586 & \textbf{0.644} & 0.511 (70\%) & 0.542 (70\%) \\
1 & 10 & 0.264 & \textbf{0.429} & 0.238 (42\%) & 0.318 (46\%) \\
1 & 20 & 0.164 & \textbf{0.286} & 0.111 (33\%) & 0.246 (36\%) \\
3 & 2 & 0.467 & \textbf{0.516} & 0.439 (63\%) & 0.451 (64\%) \\
3 & 10 & 0.189 & \textbf{0.366} & 0.178 (39\%) & 0.265 (43\%) \\
3 & 20 & 0.067 & \textbf{0.243} & 0.082 (32\%) & 0.183 (35\%) \\}
\newcommand{\SimTable}{1 & 2 & 0.744 / 0.586 & 0.731 / 0.644 & 0.703 / 0.511 & 0.688 / 0.542 \\
1 & 10 & 0.566 / 0.264 & 0.556 / 0.429 & 0.607 / 0.238 & 0.595 / 0.318 \\
1 & 20 & 0.405 / 0.164 & 0.396 / 0.286 & 0.502 / 0.111 & 0.490 / 0.246 \\
3 & 2 & 0.671 / 0.467 & 0.672 / 0.516 & 0.612 / 0.439 & 0.590 / 0.451 \\
3 & 10 & 0.528 / 0.189 & 0.506 / 0.366 & 0.521 / 0.178 & 0.505 / 0.265 \\
3 & 20 & 0.381 / 0.067 & 0.382 / 0.243 & 0.434 / 0.082 & 0.405 / 0.183 \\}
\newcommand{\TestTable}{Crossover (C$-$A at 20 vs 2\,\textmu s) & 1 & +0.021 & [-0.022, +0.066] & 0.338 & no \\
DD on bare (B$-$A, 20\,\textmu s) & 1 & +0.122 & [+0.094, +0.147] & 0.002 & yes \\
DD on encoded (D$-$C, 20\,\textmu s) & 1 & +0.135 & [+0.094, +0.181] & 0.002 & yes \\
Encoded beats bare (C$-$A, 20\,\textmu s) & 1 & -0.054 & [-0.089, -0.018] & 1.000 & no \\
Crossover (C$-$A at 20 vs 2\,\textmu s) & 3 & +0.043 & [-0.002, +0.089] & 0.031 & no \\
DD on bare (B$-$A, 20\,\textmu s) & 3 & +0.176 & [+0.152, +0.201] & 0.002 & yes \\
DD on encoded (D$-$C, 20\,\textmu s) & 3 & +0.102 & [+0.060, +0.142] & 0.002 & yes \\
Encoded beats bare (C$-$A, 20\,\textmu s) & 3 & +0.015 & [-0.018, +0.051] & 0.338 & no \\}
\newcommand{\QTable}{BARE & 1 & 8 & 15 & 0.984 [0.981, 0.987] & 100.0\% & 258 \\
ENC\_R & 1 & 10 & 15 & 0.980 [0.977, 0.983] & 99.2\% & 276 \\
ENC\_ALL & 1 & 18 & 42 & 0.987 [0.984, 0.989] & 95.5\% & 668 \\
ENC\_ALL\_CHK & 1 & 18 & 66 & 0.989 [0.987, 0.991] & 91.9\% & 985 \\
BARE & 3 & 8 & 22 & 0.975 [0.971, 0.978] & 100.0\% & 337 \\
ENC\_R & 3 & 10 & 22 & 0.969 [0.966, 0.973] & 99.2\% & 355 \\
ENC\_ALL & 3 & 18 & 42 & 0.984 [0.982, 0.987] & 94.6\% & 680 \\
ENC\_ALL\_CHK & 3 & 18 & 66 & 0.989 [0.986, 0.991] & 90.5\% & 999 \\}
\newcommand{\HQCminB}{1,336}
\newcommand{\HQCminC}{2,581}
\newcommand{\QbareLthree}{0.975}
\newcommand{\QchkLthree}{0.989}
\newcommand{\QchkKeep}{91}


\title{\bfseries Dynamical decoupling, not error detection, protects an idle quantum register on heavy-hex:
a preregistered code-space Hayden--Preskill test on IBM Heron r2, and a validated design for Quantinuum H2}
\author{Devin Phillip Davis\\ \small Agile Defense Systems LLC}
\date{21 November 2026 \quad (OSIRIS-XTOP-005, version 1.0)}

\begin{document}
\emergencystretch=3em
\maketitle

\begin{abstract}
\noindent
We asked whether encoding a quantum register in the $[[4,2,2]]$ error-detecting code improves
Hayden--Preskill-style recovery of entanglement on a superconducting processor, and whether
dynamical decoupling (DD) helps. The circuits use only Clifford gates that are fault-detecting on
$[[4,2,2]]$; we found scramblers that hide the reference exactly and decoders, of 5 to 9 gates, by
exhaustive search. On IBM's heavy-hex lattice, encoding all four registers needs 147--396 CZ gates
against 30--60 for the bare protocol, so we redesigned the experiment to encode only the idle reference
register, with a 5-CNOT star encoder that fits one degree-3 heavy-hex node. In a preregistered run
(Zenodo 10.5281/zenodo.24425416; \NCirc{} circuits, \QPUs\,s on \texttt{ibm\_marrakesh}), staggered
native-gate XY4 raised the two-pair Bell fidelity by 0.10--0.18 at a 20\,\textmu s idle window (4 of 4
DD tests supported, Holm $\alpha=0.01$). Error detection did not help: encoded post-selected fidelity
never beat the matching bare arm, and both encoding hypotheses failed. Simulation with IBM's published
noise model predicted the opposite on both counts. That points to idle noise outside that model
(crosstalk or low-frequency dephasing), which DD suppresses and detection only discards. On
all-to-all hardware the routing obstacle disappears. For Quantinuum H2 we give the full design (every
register encoded, verified preparation, flagged mid-circuit checks, 72 circuits), check every
H2-compiled circuit exactly, predict $F=\QchkLthree$ against $\QbareLthree$ for bare at depth 3 under an
H2-like noise model, and give the exact cost: \HQCminB{} HQC for the decisive comparison at 500 shots.
\end{abstract}

\section{Question}
Hayden--Preskill recovery~\cite{hp} transfers a reference's entanglement through a scrambler to
outgoing radiation, where a decoder recovers it~\cite{yk}. In any real device the reference register idles for
the whole protocol. XTOP-005 asks two practical questions: does error detection on that register (or
on all registers) improve recovered fidelity, and does dynamical decoupling during idle windows? We
measure the two-pair Bell fidelity $F=\tfrac1{16}\sum_{a,b} s_a s_b\langle\sigma_a\sigma_a\otimes
\sigma_b\sigma_b\rangle$ of the recovered pairs $(R_i,E_i)$ from 9 Pauli readout settings, with signs
$s$ from the exact noiseless frame.

\section{Protocol and gate set}
Registers $R,M,B,E$ hold two logical qubits each. Bell pairs $(R_i,M_i)$ and $(B_i,E_i)$ are prepared,
a scrambler $U_L$ acts on $(M,B)$, a transit window of length $\tau$ follows, a decoder $V_L$ acts on
$(B,E)$, and $(R,E)$ are read out. Every logical gate is fault-detecting on $[[4,2,2]]$
($S_X=XXXX$, $S_Z=ZZZZ$): transversal CNOT between blocks; $H^{\otimes4}$ (logical $H\otimes H$ then
SWAP); $S\,S^\dagger S^\dagger S$ (an exact logical CZ, virtual on Heron); and qubit relabellings
(logical SWAP and CNOT). Random search found scramblers at $L=1,2,3$ with
$I(R{:}M)=I(R{:}B)=0$ exactly, and breadth-first search over binary symplectic tableaux (up to
$2.7\times10^6$ states) found decoders of 5, 6 and 9 gates. All circuits pass an exact noiseless gate
(Stage~0) after compilation: $F=1$ and no shot rejected, in every arm, depth and setting.

\section{Why heavy-hex forced a redesign}
With every register encoded, transversal CNOTs need four parallel nearest-neighbour edges between
blocks, which heavy-hex lacks. Compiled for \texttt{FakeMarrakesh}, the encoded circuits need 147--160
CZ (341--396 with flagged checks) against 30--60 bare, and post-selected fidelity falls 6--8 points
\emph{below} the unprotected bare arm while discarding half the shots. Interleaving the blocks along a
16-qubit line makes the gates between blocks local, but each block's preparation then needs
$\approx132$ CZ. Heavy-hex cannot make both local.

The heavy-hex design therefore encodes only the reference, the register that idles throughout. A
5-CNOT star encoder (minimal on a star, found by BFS) maps $R$ into the code space on one degree-3 node.
Its inverse returns both stabilizer syndromes to two qubits, detecting every single-qubit error on $R$
in all 9 settings (versus 2 of 9 with final-readout checks alone). The overhead is about 13 CZ.

\section{Hardware experiment (IBM Heron r2)}
Four arms: A bare/idle; B bare with staggered XY4; C $R$ encoded/idle; D $R$ encoded with staggered XY4.
XY4 uses only native \texttt{x} and virtual \texttt{rz} ($Y=R_z(\pi/2)XR_z(-\pi/2)$), with pulses at odd
eighths of each delay on one heavy-hex sublattice and even eighths on the other. Depths $L\in\{1,3\}$ and
windows $\tau\in\{2,10,20\}$\,\textmu s give \NCirc{} circuits at 1000 shots each in one SamplerV2 job
(\texttt{\JobID}). The hotspot-avoiding layouts (bottom-decile $T_2$, top-decile readout or CZ error)
came from the live calibration at lock time and are shared by the idle and DD arms. The circuits (QPY
SHA-256 \texttt{\QPYsha\ldots}), hypotheses and analysis code were published before submission. The
bare and encoded arms necessarily occupy different qubits, which confounds H4; H1 is a difference of
differences on the same qubits.

\begin{table*}[t]\centering\small
\caption{Hardware fidelity $F$ (encoded arms post-selected; keep fraction in parentheses). Bold: best arm
per row.}
\setlength{\tabcolsep}{3pt}
\begin{tabular}{ccllll}\toprule
$L$ & $\tau$/\textmu s & A bare & B bare+DD & C enc $R$ & D enc $R$+DD\\\midrule
\HWTable
\bottomrule\end{tabular}\end{table*}

\begin{table*}[t]\centering\small
\caption{Preregistered one-sided tests (8 total, Holm-corrected, 4000-resample bootstrap).}
\setlength{\tabcolsep}{5pt}
\begin{tabular}{lcrlrl}\toprule
Test & $L$ & Est. & 99\% CI & $p_\text{Holm}$ & Supp.\\\midrule
\TestTable
\bottomrule\end{tabular}\end{table*}

\paragraph{Results.} DD was the largest effect in the experiment: at 20\,\textmu s it lifted bare fidelity
by $0.12$ ($L=1$) and $0.18$ ($L=3$), and lifted encoded fidelity by $0.14$ and $0.10$. All four DD tests
were supported. Error detection did not help. Encoded post-selected fidelity was never above the
matching bare arm, the keep fraction fell to 32--36\% at 20\,\textmu s, and neither the crossover (H1) nor
the long-idle advantage (H4) was supported; at $L=1$ the encoded arm was significantly worse.

\paragraph{The simulator got both effects backwards.} Table~\ref{tab:sim} compares the preregistered
Stage-1 predictions (Aer with \texttt{FakeMarrakesh} noise) with hardware. The simulator predicted no DD
benefit (its noise has $T_1$, $T_2$ and gate error but no ZZ crosstalk or low-frequency noise) and an
encoding win at long idles. On hardware, bare fidelity at 10\,\textmu s was 0.26 against a predicted
0.57, and DD recovered much of the gap. Most idle error on this device is therefore coherent or
correlated noise outside the published model. DD refocuses it; a distance-2 detection code can only
discard the affected shots, and here it discarded most of them.

\begin{table}[t]\centering\footnotesize
\caption{Stage-1 prediction / hardware, same cells as Table 1.}\label{tab:sim}
\setlength{\tabcolsep}{3pt}
\begin{tabular}{ccllll}\toprule
$L$ & $\tau$ & A & B & C & D\\\midrule
\SimTable
\bottomrule\end{tabular}\end{table}

\section{Quantinuum H2 design}
All-to-all connectivity removes the obstacle of Section~3: every transversal CNOT is one native
two-qubit gate, and mid-circuit measurement with qubit reuse makes flagged checks cheap. The H2 design
compares four protection levels on the same logical program ($L\in\{1,3\}$, 9 settings, 72 circuits):
Q0 bare (8 qubits); Q1 $R$ encoded (the IBM arm, for a cross-platform comparison; 10 qubits); Q2 all
registers encoded with verified $|00\rangle_L/|{+}{+}\rangle_L$ preparation (17--18 qubits); and Q3, Q2
plus flagged weight-4 $X$ and $Z$ checks on $R$ and $E$ after decoding (18 qubits). DD is dropped
because idling on H2 is not a user-timed delay.

Every circuit was compiled offline for H2-1 with \texttt{pytket-quantinuum} and the compiled circuits
were verified exactly ($F=1$, no shot rejected; all 72). Under an H2-like Pauli noise model
(published magnitudes: $3\times10^{-5}$ one-qubit, $10^{-3}$ two-qubit, $10^{-3}$ SPAM; not Quantinuum's
emulator, and with no memory error), full encoding with checks gives the highest fidelity at both depths
(Table~\ref{tab:q}). At $L=3$, $F=\QchkLthree$ against $\QbareLthree$ bare, keeping \QchkKeep\% of shots, with
non-overlapping 99\% intervals. This is the regime IBM's heavy-hex could not reach.

\begin{table*}[t]\centering\small
\caption{H2 design: H2-1-compiled two-qubit gates, predicted $F$ (99\% CI), keep fraction, and HQC for 9
settings at 500 shots ($\mathrm{HQC}=5+C(N_{1q}+10N_{2q}+5N_m)/5000$).}\label{tab:q}
\setlength{\tabcolsep}{2pt}
\begin{tabular}{lcccllr}\toprule
Arm & $L$ & qubits & $N_{2q}$ & $F$ (pred.) & keep & HQC\\\midrule
\QTable
\bottomrule\end{tabular}\end{table*}

\paragraph{Cost.} The full 72-circuit matrix costs \HQCtotalA{} HQC at 200 shots, \HQCtotalB{} at 500
and \HQCtotalC{} at 1000. The decisive comparison alone (Q0 against Q3 at $L=3$, 9 settings) costs
\HQCminB{} HQC at 500 shots, enough to resolve the predicted 0.014 difference, or \HQCminC{} at 1000.
The compiled H2-1 circuits are included (pytket JSON, SHA-256 \texttt{\QSha\ldots}) so the exact
submission can be inspected.

\section{Limitations}
The heavy-hex experiment ran once, on one device, at one calibration. The encoded and bare arms use
different physical qubits. The H2 predictions rest on an approximate Pauli noise model without memory
or transport errors and have not been tested on Quantinuum hardware or its emulator. The scramblers
are Clifford circuits, so this tests entanglement transfer through a transversal-gate scrambler, not
scrambling in the sense of black-hole dynamics.

\section{Data and code}
Preregistration, code, lock, QPY circuits and Stage-1 predictions: 10.5281/zenodo.24425416. This record
adds raw counts, the analysis output with exec-canon-1 receipts, the Quantinuum design module and the
compiled H2-1 circuits. Source: \texttt{xtop005/} in \texttt{quantum-osiris/osiris-cli}.

\begin{thebibliography}{9}\small
\bibitem{hp} P.~Hayden and J.~Preskill, ``Black holes as mirrors,'' JHEP 09 (2007) 120.
\bibitem{yk} B.~Yoshida and A.~Kitaev, ``Efficient decoding for the Hayden--Preskill protocol,''
arXiv:1710.03363.
\bibitem{prereg} D.~P.~Davis, Preregistration: OSIRIS-XTOP-005, 10.5281/zenodo.24425416 (2026).
\bibitem{ro} D.~P.~Davis, TESSERACT-RO-06/07, 10.5281/zenodo.24420902 (2026).
\end{thebibliography}
\end{document}
